\section{Fundamentação Teórica}
\label{sec:fundamentacao}

\subsection{Programação Linear e Otimização Inteira}
\label{sec:prog-linear}

A otimização matemática trata da escolha, dentre um conjunto de alternativas viáveis, daquela que minimiza (ou maximiza) uma medida de desempenho. Quando tanto a função objetivo quanto as restrições que delimitam as alternativas são expressas por relações lineares, o problema é dito de \emph{programação linear}. Esta seção introduz os conceitos de programação linear, sua extensão ao domínio inteiro e o caso misto, que constitui a base formal do modelo proposto neste trabalho.

\paragraph{Programação Linear (PL).}
Um problema de programação linear busca otimizar uma função objetivo linear sujeita a um conjunto de restrições também lineares. Na forma padrão, pode ser escrito como
\begin{align}
    \min_{x}\quad & c^{\top} x \notag \\
    \text{sujeito a}\quad & A x \le b, \label{eq:pl-padrao} \\
    & x \ge 0, \notag
\end{align}
onde $x \in \mathbb{R}^n$ é o vetor de variáveis de decisão (contínuas), $c \in \mathbb{R}^n$ define os coeficientes da função objetivo, e $A \in \mathbb{R}^{m \times n}$ e $b \in \mathbb{R}^m$ descrevem as restrições. O conjunto de pontos que satisfazem todas as restrições é a \emph{região factível}, que forma um poliedro convexo. Um resultado central da otimização linear estabelece que, se o problema admite solução ótima, então existe um \emph{vértice} (ponto extremo) da região factível que é ótimo \cite[Propriedade~2.2]{arenales2007pesquisa}; de modo equivalente, basta procurar o ótimo entre as soluções básicas factíveis, que correspondem aos vértices da região. Essa propriedade fundamenta o método \emph{simplex}, que percorre vértices adjacentes da região factível, melhorando a função objetivo a cada passo até alcançar o ótimo \cite{arenales2007pesquisa}.

\paragraph{Programação Linear Inteira (PLI).}
Em muitas aplicações práticas, as variáveis de decisão representam grandezas indivisíveis ou escolhas lógicas e, portanto, devem assumir apenas valores inteiros. Acrescentar a exigência $x \in \mathbb{Z}^n$ ao problema~\eqref{eq:pl-padrao} caracteriza a \emph{programação linear inteira}. Um caso particularmente importante é o das variáveis \emph{binárias} ($x \in \{0,1\}^n$), capazes de modelar decisões do tipo sim/não, como ``visitar ou não um nó'' ou ``ativar ou não uma rota''. Diferentemente da PL, a PLI é, em geral, computacionalmente muito mais difícil de resolver: a restrição de integralidade torna o conjunto factível discreto (não convexo) e o número de combinações inteiras a examinar cresce rapidamente com a quantidade de variáveis, tornando inviável a enumeração completa das soluções em instâncias de grande porte \cite{arenales2007pesquisa}.

\paragraph{Programação Linear Inteira Mista (MILP).}
Quando apenas parte das variáveis é restrita a valores inteiros, enquanto as demais permanecem contínuas, tem-se um problema de \emph{programação linear inteira mista} (MILP, do inglês \textit{Mixed-Integer Linear Programming}) \cite{arenales2007pesquisa}. Essa formulação é especialmente expressiva para problemas de roteamento e escalonamento, nos quais variáveis binárias capturam decisões discretas (por exemplo, qual aresta percorrer) e variáveis contínuas representam grandezas físicas acumuladas (por exemplo, energia ou tempo).

\paragraph{Resolução por \textit{branch-and-bound}.}
A principal técnica para resolver problemas de PLI e MILP de forma exata é o \emph{branch-and-bound} \cite{land1960branchbound, arenales2007pesquisa}. A ideia central é dividir o problema em subproblemas menores (\emph{branching}) e usar limitantes para descartar (\emph{bounding}) aqueles que comprovadamente não podem conter a solução ótima, evitando testar todas as combinações possíveis. Esse método, em conjunto com técnicas de planos de corte, está implementado em resolvedores de otimização como o Gurobi, adotado neste trabalho.

\paragraph{Relação com este trabalho.}
O problema de planejamento de voo abordado neste trabalho é formulado como um MILP: variáveis binárias descrevem a posição do UAV, os deslocamentos e as coletas, enquanto variáveis contínuas registram a energia acumulada e a \emph{Age of Information} dos sensores. Termos não lineares, como o produto entre a AoI e a decisão de coleta, são convertidos em restrições lineares por meio de técnicas de linearização (Seção~\ref{sec:modelagem}, Equações~\eqref{eq:lin-a}--\eqref{eq:lin-c}), preservando a estrutura linear necessária à resolução exata pelo Gurobi.

\subsection{IoT e WSN}

A Internet das Coisas (IoT, do inglês \textit{Internet of Things}) é um paradigma que integra objetos físicos (sensores, atuadores e dispositivos embarcados) à infraestrutura digital por meio de protocolos de comunicação sem fio, permitindo a coleta, transmissão e processamento autônomo de dados do ambiente \cite{yang2020surveyUAVIoT}. Dentro desse paradigma, as redes de sensores sem fio (WSNs, do inglês \textit{Wireless Sensor Networks}) constituem uma das tecnologias de suporte mais consolidadas: consistem em conjuntos de nós de baixo consumo energético, distribuídos em uma área de interesse, responsáveis por amostrar grandezas físicas (temperatura, umidade, vibração, entre outras) e encaminhá-las periodicamente a um ponto de coleta centralizado \cite{liu2018uavWSN}.

Em aplicações de monitoramento ambiental, agricultura de precisão, vigilância urbana e resposta a desastres, os nós sensores frequentemente operam em áreas de difícil acesso ou sem cobertura de infraestrutura fixa. Nesse contexto, UAVs têm sido explorados como \emph{sinks} móveis, percorrendo a área de missão para coletar dados diretamente dos sensores \cite{wei2022survey,messaoudi2023survey}.

\subsection{Age of Information (AoI)}

Uma métrica central para avaliar a qualidade da informação coletada nessas redes é a \emph{Age of Information} (AoI), introduzida para quantificar a \emph{frescor} dos dados disponíveis em um ponto de destino \cite{messaoudi2023survey}. Intuitivamente, a AoI mede há quanto tempo a última atualização de um determinado sensor foi recebida. Quanto maior a AoI, mais desatualizada é a informação disponível.

\paragraph{Definição formal.}

No modelo de tempo discreto adotado neste trabalho, o horizonte de missão é dividido em \emph{slots} de duração $\Delta t$. Para cada sensor $j \in S$, define-se $A_{j,t}$ como a \textbf{idade da informação} do sensor $j$ no início do slot $t$, medida em número de slots desde a última coleta. Sua dinâmica é dada por:

\begin{align}
    A_{j,1} &= A_j^0, \qquad \forall j \in S, \tag{AoI-0} \\[4pt]
    A_{j,t+1} &=
    \begin{cases}
        0, & \text{se } v_{j,t} = 1 \text{ (coleta realizada no slot } t\text{)}, \\[4pt]
        A_{j,t} + 1, & \text{se } v_{j,t} = 0 \text{ (sem coleta)},
    \end{cases}
    \tag{AoI-1}
\end{align}

onde $A_j^0$ é o valor inicial carregado do estado persistido entre missões e $v_{j,t} \in \{0,1\}$ é a variável de decisão que indica se o sensor $j$ é coletado no slot $t$.

Quando o UAV coleta o sensor $j$ no slot $t$, a AoI é zerada no slot seguinte, pois o destino passa a dispor de uma atualização recente. Caso contrário, a AoI cresce em uma unidade por slot, refletindo o envelhecimento da informação. Maximizar a soma $\sum_{j,t} A_{j,t}\, v_{j,t}$, isto é, priorizar sensores com AoI mais alta no momento da coleta, é equivalente a maximizar o \emph{ganho informacional} acumulado pela missão, incentivando visitas a nós cujos dados estejam mais desatualizados.

\subsection{UAVs}

Veículos aéreos não tripulados (UAVs, do inglês \textit{Unmanned Aerial Vehicles}), popularmente chamados de drones, são aeronaves controladas remotamente ou de forma autônoma, sem piloto a bordo \cite{messaoudi2023survey}. Originalmente desenvolvidos para fins militares, os UAVs passaram a ser amplamente utilizados em aplicações civis como mapeamento, agricultura de precisão, vigilância, entrega de cargas e, em particular, coleta de dados em redes IoT \cite{wei2022survey}.

Do ponto de vista da propulsão, os UAVs mais relevantes para coleta de dados em WSNs se dividem em dois grandes grupos:

\begin{itemize}
    \item \textbf{Asa fixa (\textit{fixed-wing})}: aeronaves que geram sustentação pelo perfil aerodinâmico das asas durante o deslocamento horizontal. Possuem maior autonomia e velocidade de cruzeiro, mas não conseguem pairar estáticos sobre um ponto e exigem velocidade mínima de voo para manter sustentação.
    \item \textbf{Asa rotativa (\textit{rotary-wing})}: aeronaves com um ou mais rotores horizontais (helicópteros, quadricópteros, hexacópteros). Têm menor autonomia e velocidade máxima, mas são capazes de realizar \emph{hover} (permanecer estacionários no ar), o que facilita a comunicação direta com sensores fixos no solo \cite{zeng2019rotarywing}.
\end{itemize}

Neste trabalho, adota-se o modelo de asa rotativa, coerente com os UAVs utilizados nos experimentos (DJI Matrice 300 RTK e Parrot Anafi USA).

\subsubsection{Modos de coleta}
\label{sec:modos-coleta}

A estratégia de sobrevoo adotada pelo UAV determina como e quando a comunicação com os sensores ocorre. Três modos principais são identificados na literatura \cite{messaoudi2023survey,wei2022survey}:

\begin{itemize}
    \item \textbf{\textit{Hovering}}: o UAV desloca-se até a posição do sensor e paira estático acima dele pelo tempo necessário para completar a transmissão. Garante máxima qualidade de enlace (menor distância e velocidade relativa nula), à custa de maior consumo energético para manter o voo estacionário.
    \item \textbf{\textit{Fly-by}}: o drone sobrevoa o sensor em movimento contínuo, realizando a coleta em trânsito. Reduz o tempo total de missão e o consumo associado ao hover, porém a janela de comunicação é curta e a qualidade do enlace pode ser comprometida pela variação da distância e pelo efeito Doppler.
    \item \textbf{Híbrido}: combina os dois modos anteriores, selecionando a estratégia mais adequada para cada sensor conforme seus requisitos de enlace e localização.
\end{itemize}

No modelo adotado neste trabalho, a coleta ocorre exclusivamente em modo \textit{hovering}: o UAV deve permanecer parado sobre o sensor durante um slot inteiro para que a transmissão seja considerada bem-sucedida. Esse pressuposto é modelado pela restrição $v_{j,t} = x_{jjt}$, que vincula a coleta ao \emph{self-loop} do grafo de movimento (ver Seção~\ref{sec:modelagem}).

\subsubsection{Restrições operacionais}

A operação de um UAV de asa rotativa é condicionada por dois limites principais:

\begin{itemize}
    \item \textbf{Capacidade de bateria ($E^{\max}$)}: a energia armazenada na bateria limita o tempo total de missão. O consumo depende da velocidade de deslocamento e do tempo em hover, conforme o modelo de potência de Zeng et al.~\cite{zeng2019rotarywing} descrito na Seção~\ref{sec:modelagem}.
    \item \textbf{Horizonte temporal ($T_f$)}: o número máximo de slots disponíveis para a missão, que define o comprimento do plano de voo a ser otimizado.
\end{itemize}

A interação entre esses dois limites é central para o problema de otimização formulado neste trabalho: aumentar o número de sensores visitados exige mais deslocamentos e mais tempo em hover, elevando o consumo e eventualmente inviabilizando a missão dentro do orçamento energético disponível.

\subsection{Modelo de energia de UAV de asa rotativa}

O consumo energético de UAVs de asa rotativa é fortemente dependente do regime de voo: manter-se estacionário (\emph{hover}) exige potência elevada para sustentar o peso do veículo contra a gravidade, enquanto o deslocamento horizontal apresenta um ponto ótimo de eficiência situado entre o custo de sustentação e o arrasto aerodinâmico.

Para modelar esse comportamento com fidelidade física, este trabalho adota o modelo de potência proposto por Zeng et al.~\cite{zeng2019rotarywing}, que decompõe a potência instantânea de um rotor em três componentes:

\begin{enumerate}
    \item \textbf{Potência de perfil das pás}: cobre o arrasto viscoso das pás ao girarem. Em translação horizontal, as pás avançam e recuam assimetricamente contra o fluxo de ar, fazendo esta parcela crescer proporcionalmente ao quadrado da velocidade.
    \item \textbf{Potência induzida}: energia transferida ao ar para gerar sustentação (empuxo). Domina em \emph{hover}, decresce até um mínimo com o aumento da velocidade e volta a crescer em velocidades elevadas.
    \item \textbf{Potência de arrasto parasita}: captura o arrasto aerodinâmico da fuselagem e de componentes expostos ao fluxo; cresce proporcionalmente ao cubo da velocidade e é dominante apenas em altas velocidades.
\end{enumerate}

A expressão matemática da potência $P(V)$ em função da velocidade horizontal $V$ é apresentada na Seção~\ref{sec:modelagem} (Equação~\eqref{eq:potencia}). Em \emph{hover} ($V = 0$), o modelo simplifica-se para $P_{\mathrm{hov}} = P_0 + P_i$, onde $P_0$ é a potência de perfil e $P_i$ é a potência induzida, ambas avaliadas em velocidade nula. Com os parâmetros do DJI Matrice 300 RTK ($P_0 = 79{,}86$\,W, $P_i = 88{,}63$\,W), obtém-se $P_{\mathrm{hov}} \approx 168{,}5$\,W.

Os parâmetros físicos completos utilizados nos experimentos, obtidos de \cite{zeng2019rotarywing}, são: $P_0 = 79{,}86$\,W, $P_i = 88{,}63$\,W, $U_{\mathrm{tip}} = 120$\,m/s (velocidade da ponta da pá), $v_0 = 4{,}03$\,m/s (velocidade induzida em \emph{hover}), $\rho = 1{,}225$\,kg/m³ (densidade do ar) e $C_{DS} = 0{,}01509$\,m² (coeficiente de arrasto parasita efetivo, dado por $c_d \cdot s \cdot A$, com $c_d = 0{,}6$, $s = 0{,}05$ e $A = 0{,}503$\,m²). Esses parâmetros são adotados nos dois conjuntos de experimentos; a diferença entre os UAVs está na capacidade de bateria $E^{\max}$: $2$\,MJ para o DJI Matrice 300 RTK e $141{,}372$\,kJ para o Parrot Anafi USA (capacidade nominal real do equipamento).

O custo energético de cada transição no grafo de movimento é obtido multiplicando a potência calculada pela duração do slot $\Delta t = 10$\,s (Equação~\eqref{eq:custo-energia}, Seção~\ref{sec:modelagem}). As restrições de acumulação e capacidade de energia no modelo MILP são apresentadas na Seção~\ref{sec:modelagem} (Equações~\eqref{eq:energia-inicial}--\eqref{eq:energia-capacidade}).

\subsection{Modelo de comunicação}

Em sistemas UAV-IoT, modelos de comunicação detalhados descrevem explicitamente o canal de rádio (perda de percurso, desvanecimento, relação sinal-ruído e protocolos de controle de acesso ao meio) condicionando a taxa de sucesso à qualidade do enlace~\cite{wei2022survey,messaoudi2023survey}. Esses modelos são relevantes quando a mobilidade do UAV provoca variações significativas de distância ou ângulo durante a coleta.

Neste trabalho, adota-se um modelo de comunicação binário simplificado, coerente com o modo de coleta \emph{hovering} descrito na Seção~\ref{sec:modos-coleta}. O pressuposto central é que a transmissão de dados ocorre com sucesso garantido se, e somente se, o UAV permanecer estacionário diretamente acima do sensor durante um slot completo de duração $\Delta t$.

Não são modelados aspectos de camada de enlace como taxa de dados, modulação, interferência entre sensores ou protocolo MAC. A coleta é representada pela variável binária $v_{j,t} \in \{0,1\}$, vinculada ao \emph{self-loop} de permanência do UAV sobre o nó sensor (Equação~\eqref{eq:coleta}, Seção~\ref{sec:modelagem}).

Esta simplificação é adequada por três motivos:

\begin{enumerate}
    \item \textbf{Distância mínima:} ao pairar sobre o sensor, o UAV minimiza a distância de enlace, garantindo SNR suficiente para tecnologias de curto alcance como IEEE 802.15.4, Zigbee ou LoRa.
    \item \textbf{Canal estacionário:} em \emph{hover}, a velocidade relativa entre drone e sensor é nula, eliminando o efeito Doppler e reduzindo as flutuações de desvanecimento.
    \item \textbf{Foco em roteamento:} o objetivo da formulação é otimizar a trajetória e o escalonamento de visitas; o desempenho do enlace de rádio, condicionalmente à presença de \emph{hover}, é tratado como garantido. Modelos de canal mais detalhados podem ser incorporados em extensões futuras.
\end{enumerate}
